\chapter{Occupancy Prediction}
\label{chap:occupancy-prediction}

Occupancy prediction is used to estimate whether a hotel room is
\textit{Occupied}, \textit{Vacant}, or under \textit{Cleaning}. This
information supports later lighting and HVAC decisions. For example, a vacant
room can use more energy-efficient settings, whereas an occupied room should
give greater priority to guest comfort.

\section{Input Data and Features}

The occupancy models use PIR motion records together with reservation and guest
information. Important inputs include recent motion activity, number of adults
and children, active reservation status, room type, hour of day, and day of
week. Time values are represented cyclically so that neighboring periods, such
as 23:00 and 00:00, remain close to each other.

The prediction target is the room state one hour in the future. Since sensor
records are sampled every five minutes, the model uses 12 samples to represent
one hour.

\section{Implemented Models}

Two occupancy models are available. The Random Forest model uses engineered
tabular features, including recent motion statistics and reservation context.
The Transformer model processes a sequence of recent room observations and
learns temporal relationships between them.

The Flask endpoint \texttt{/api/predict\_occupancy} returns the predicted
class, confidence score, class probabilities, and diagnostic values such as
motion rate. The user interface allows the Random Forest or Transformer model
to be selected.

\section{Performance}

The Random Forest achieved 98.9\% accuracy, while the Transformer achieved
98.8\%. Both models performed strongly for Occupied and Vacant rooms.
Cleaning was more difficult because it has considerably fewer examples in the
dataset.


\chapter{Guest Persona Classification}
\label{chap:persona-classification}

Guest persona classification identifies repeated lighting and temperature
preferences. The predicted persona is used as contextual information by the
recommendation systems.

\section{Lighting Personas}

The lighting personas used by the application are:

\begin{itemize}
    \item \textit{Balanced}: prefers moderate lighting that changes with
    activity and time;
    \item \textit{Routine}: follows repeated daily lighting patterns;
    \item \textit{StaticBright}: generally prefers high brightness;
    \item \textit{StaticDim}: generally prefers low brightness;
    \item \textit{NightFocused}: uses lighting mainly during evening and night;
    \item \textit{Housekeeping}: uses bright task-oriented lighting; and
    \item \textit{Unknown}: does not match a known pattern or has insufficient
    history.
\end{itemize}

Lighting persona prediction uses lamp levels, time-of-day patterns, PIR motion,
occupant counts, and the proportion of time that lamps remain active. A Random
Forest model uses daily summary features, while a Transformer analyzes a
complete sequence of five-minute lighting observations.

\section{Temperature Personas}

The temperature personas are AlwaysOnComfort, EnergySaver, Reactive,
Preconditioning, and Housekeeping. They represent different relationships
between comfort, HVAC operation, room conditions, and energy use.

The temperature persona model uses recent room temperature, outside
temperature, setpoint, HVAC mode, occupancy state, room characteristics, and
PIR activity. It is implemented as a Transformer classifier.

\section{Performance}

The Random Forest lighting persona model achieved 88.5\% accuracy, while the
lighting Transformer achieved 92.5\%. The temperature persona Transformer
achieved 80.55\% accuracy.


\chapter{Lighting Recommendation System}
\label{chap:lighting-recommendation}

The lighting recommendation system generates a brightness level for each lamp
in a selected room. It combines predicted occupancy, lighting persona, current
lamp state, reservation status, PIR activity, and recent lighting history.

\section{Recommendation Process}

The workflow consists of the following steps:

\begin{enumerate}
    \item The user selects a room, timestamp, history window, and model.
    \item Occupancy and lighting persona are predicted.
    \item Recent lighting measurements are retrieved from the database.
    \item A recommended level is calculated for every lamp.
    \item Hardware constraints are applied.
    \item Current and recommended energy values are compared.
\end{enumerate}

The system distinguishes between dimmable LEDs and non-dimmable bulbs.
Dimmable lamps may receive levels between 0 and 80, while non-dimmable lamps
are restricted to off or full level. Lamps in vacant rooms are normally turned
off, while Cleaning conditions may require brighter lighting.

\section{Recommendation Models}

The application supports a HistGradientBoosting regressor and a Transformer
regressor. The first model uses current values and historical summary
features. The Transformer processes a sequence of recent five-minute records
and predicts the level for the next five-minute interval.

The result includes the current and recommended level for each lamp, the model
used, an explanation, and estimated energy savings. These values are displayed
in the React dashboard as comparison bars and summary cards.


\chapter{Temperature and HVAC Recommendation}
\label{chap:temperature-recommendation}

The HVAC recommendation system suggests a suitable room-temperature setpoint.
Its goal is to reduce unnecessary energy consumption without ignoring guest
comfort.

\section{Recommendation Inputs}

The model uses recent room temperature, outside temperature, current setpoint,
ideal temperature, HVAC mode, room size, PIR motion, occupancy prediction, and
temperature persona prediction. The input is represented as a sequence of 24
five-minute observations, corresponding to two hours of recent history.

\section{Setpoint Recommendation}

The Transformer recommends a setpoint between
$16\,^{\circ}\mathrm{C}$ and $28\,^{\circ}\mathrm{C}$. The result is rounded
to the nearest $0.5\,^{\circ}\mathrm{C}$. Vacant rooms are allowed to use
wider energy-saving settings, while the AlwaysOnComfort persona receives
settings closer to the preferred comfort temperature.

The API returns the current and recommended setpoints, the required action
(\textit{raise}, \textit{lower}, or \textit{keep}), the expected HVAC mode,
and an immediate energy estimate.

\section{Performance}

The temperature recommendation Transformer achieved a mean absolute error of
$0.217\,^{\circ}\mathrm{C}$, an RMSE of
$0.310\,^{\circ}\mathrm{C}$, and an $R^2$ score of 0.962.


\chapter{Energy Consumption Analysis}
\label{chap:energy-analysis}

\section{Lighting Energy}

Lighting power is calculated from voltage and current measurements stored in
the dimmer reference file:

\[
    P = V \times I.
\]

Each lighting record represents a five-minute interval. Therefore, the energy
of one sample is:

\[
    E_{\mathrm{sample}} = P \times \frac{5}{60}.
\]

The system compares actual lighting consumption with maximum-brightness and
AI-recommended scenarios. It reports total energy, per-lamp consumption,
dimmable-lamp savings, and percentage savings.

\section{HVAC Energy}

Direct HVAC meter readings are not available, so power is estimated using HVAC
mode, room temperature, setpoint, outside temperature, room size, occupancy,
and PIR activity. Heating and cooling use different rated-power limits.

The estimated power of each sample is converted to energy using the same
five-minute interval. Actual HVAC use is compared with a baseline in which the
dominant heating or cooling mode operates continuously at rated power.

The dashboard presents total HVAC energy, heating and cooling contributions,
active operating time, baseline consumption, and estimated savings. HVAC
results should be interpreted as estimates rather than direct measurements.


\chapter{Machine Learning Models}
\label{chap:machine-learning-models}

The project combines classical machine-learning models and Transformer neural
networks.

\section{Classical Models}

Random Forest classifiers are used for occupancy and lighting persona
prediction. They are suitable for mixed numerical and categorical features and
can model nonlinear relationships. HistGradientBoosting is used for lighting
recommendation because it performs efficiently on large tabular datasets.

\section{Transformer Models}

Transformers are used for occupancy prediction, persona classification,
lighting recommendation, and temperature recommendation. They process
time-ordered sensor observations and learn relationships between different
parts of a sequence.

Classification models return class probabilities, while regression models
return a lighting level or temperature setpoint. The trained models and their
metadata are loaded by the Flask backend when a prediction is requested.

\section{Evaluation Metrics}

Classification performance is measured with accuracy, precision, recall,
F1-score, and confusion matrices. Regression performance is measured with
mean absolute error, root mean squared error, and $R^2$.

Chronological splitting is used for most time-series tasks so that earlier
observations are used for training and later observations are used for
evaluation.


\chapter{Results and Evaluation}
\label{chap:results}

\section{Model Results}

\begin{table}[htbp]
\centering
\caption{Summary of model evaluation results}
\label{tab:project-model-results}
\begin{tabular}{|p{4.0cm}|p{3.6cm}|p{5.0cm}|}
\hline
\textbf{Task} & \textbf{Model} & \textbf{Result} \\ \hline
Occupancy prediction & Random Forest &
Accuracy: 98.9\% \\ \hline
Occupancy prediction & Transformer &
Accuracy: 98.8\% \\ \hline
Lighting persona & Random Forest &
Accuracy: 88.5\% \\ \hline
Lighting persona & Transformer &
Accuracy: 92.5\% \\ \hline
Temperature persona & Transformer &
Accuracy: 80.55\% \\ \hline
Lighting recommendation & HistGradientBoosting &
MAE: 0.105; $R^2$: 0.999 \\ \hline
Lighting recommendation & Transformer &
MAE: 1.027; $R^2$: 0.985 \\ \hline
Temperature recommendation & Transformer &
MAE: $0.217\,^{\circ}\mathrm{C}$; $R^2$: 0.962 \\ \hline
\end{tabular}
\end{table}

\section{Discussion}

The results show that the application can predict the main occupancy states
and lighting personas with high accuracy. The Transformer improved lighting
persona classification by learning daily temporal patterns. The occupancy
models performed well overall, although the Cleaning class remained difficult
because it was underrepresented.

Both recommendation models reproduced their efficient target values with low
error. The lighting models estimated a reduction of approximately 27--29\% in
the brightness-level energy proxy. However, these targets are based on the
project's energy-saving policies and should be validated in a real hotel.


\chapter{Limitations and Future Improvements}
\label{chap:limitations}

The project has several limitations. PIR sensors measure movement rather than
direct presence, so an inactive guest may appear absent. The Cleaning class
contains relatively few observations, which reduces prediction quality for
that state.

Lighting and temperature recommendation targets are created from predefined
comfort and efficiency policies. High model scores therefore show that the
models learned these policies, but they do not prove that the policies are
optimal for every guest.

HVAC energy is estimated rather than measured with a dedicated meter. The
model does not fully represent factors such as humidity, open windows, solar
gain, insulation, or equipment efficiency.

Future improvements could include:

\begin{itemize}
    \item collecting real hotel data over longer periods;
    \item adding direct lighting and HVAC energy meters;
    \item collecting guest comfort feedback;
    \item improving minority-class occupancy prediction;
    \item monitoring prediction confidence and data drift;
    \item adding authentication and privacy controls; and
    \item connecting recommendations to real-time IoT controllers.
\end{itemize}


\chapter{Conclusion}
\label{chap:conclusion}

This project presents an AI-driven system for smart hotel-room monitoring,
personalization, and energy analysis. It combines hotel reservations, PIR
motion, lighting measurements, temperature records, weather information, and
electrical dimmer data.

The application predicts room occupancy and guest personas before generating
lighting and HVAC recommendations. It also compares actual, maximum, and
recommended energy scenarios through an interactive React and Flask
dashboard.

The experimental results demonstrate strong performance for occupancy
prediction, lighting persona classification, and recommendation tasks. The
system shows how sensor data and machine learning can support more comfortable
and energy-efficient hotel operation. Real-world measurements and guest
feedback are still required before full deployment.
